\documentclass[conference]{IEEEtran}

\usepackage{graphicx}
\usepackage{amsmath,amssymb}
\usepackage{float}
\usepackage{booktabs}
\usepackage{multirow}

\begin{document}

\title{
    \textbf{Benchmark Report: 538.imagick\_r} \\
    \text{Register Spill Analysis on RISC-V via gem5}
}

\author{
    \IEEEauthorblockN{Lütfullah Burak Kaya}
    \IEEEauthorblockA{
        Ozyegin University \\
        burak.kaya.50711@ozu.edu.tr
    }
    \and
    \IEEEauthorblockN{Ismail Akturk}
    \IEEEauthorblockA{
        Ozyegin University \\
        ismail.akturk@ozyegin.edu.tr
    }
}

\newcommand{\LongDate}{%
    \number\day\space
    \ifcase\month
    \or January\or February\or March\or April\or May\or June%
    \or July\or August\or September\or October\or November\or December%
    \fi
    \space \number\year
}

\twocolumn[
\begin{@twocolumnfalse}
\maketitle
\begin{center}{\small \LongDate}\end{center}
\vspace{1em}
\end{@twocolumnfalse}
]

% ================================================

\begin{abstract}
This report presents the register spill characterization of the
\texttt{538.imagick\_r} benchmark from SPEC CPU2017, executed on a
RISC-V (RV64GC) target using the gem5 TimingSimpleCPU simulator.
A custom SpillDetector was integrated into gem5 to count store-then-load
patterns on the stack within a user-defined Region of Interest (ROI).
The test dataset (\texttt{test\_input.tga}, with shear, resize, negate,
and alpha operations) was used as the input.
Results show a spill rate of \textbf{2.06\%} within the ROI,
corresponding to 2.4 million detected spill events across 117 million
ROI instructions. The benchmark completed in only 93 seconds of host
time, making it the fastest simulation in this study. Despite the
relatively modest spill rate, the large ImageMagick codebase exhibits
non-trivial register pressure across its image transformation pipeline.
\end{abstract}

% =========================================================
\section{Benchmark Overview}

\texttt{538.imagick\_r} is the SPEC CPU2017 rate variant of
ImageMagick's \texttt{convert} utility, a widely used open-source
image processing tool. The benchmark applies a sequence of image
transformations to an input TARGA (\texttt{.tga}) image: geometric
operations (shear, resize), pixel-level operations (negate), and
channel masking (alpha removal). The workload exercises a large,
diverse codebase covering pixel format conversion, resampling filters,
geometric transformation matrices, and color space arithmetic.

The RISC-V binary (\texttt{imagick\_r.riscv}) was compiled with gem5
ROI markers inserted in \texttt{src/utilities/convert.c}, wrapping
the entire \texttt{ConvertMain()} call from \texttt{main()}:

\begin{itemize}
    \item \texttt{m5\_work\_begin(0,0)} --- before \texttt{ConvertMain()}
    \item \texttt{m5\_work\_end(0,0)}   --- after  \texttt{ConvertMain()}
\end{itemize}

\texttt{ConvertMain()} parses the command-line options and dispatches
all image transformation operations, making it the sole entry point
for all benchmark computation.

% =========================================================
\section{Simulation Setup}

\subsection{Input Datasets}

The test and refrate datasets differ in image size, input image, and
the set of applied transformations (Table~\ref{tab:inputs}).
The refrate dataset applies more computationally expensive operations
(edge detection, resampling, emboss, color space conversion,
mean-shift filtering).

\begin{table}[H]
\centering
\caption{imagick\_r Input Dataset Sizes}
\label{tab:inputs}
\begin{tabular}{lp{5cm}}
\toprule
\textbf{Dataset} & \textbf{Operations} \\
\midrule
test    & \texttt{-shear 25 -resize 640x480 -negate -alpha Off} \\
refrate & \texttt{-edge 41 -resample 181\% -emboss 31 -colorspace YUV -mean-shift 19x19+15\% -resize 30\%} \\
\bottomrule
\end{tabular}
\end{table}

The \textbf{test} dataset was selected for this run. The input image
is \texttt{test\_input.tga}; the operations applied are:
\begin{enumerate}
    \item \textbf{Shear} by 25 degrees --- affine transform with background fill
    \item \textbf{Resize} to $640 \times 480$ --- Lanczos resampling
    \item \textbf{Negate} --- per-pixel channel inversion
    \item \textbf{Alpha Off} --- remove alpha channel
\end{enumerate}
The result is written to \texttt{test\_output.tga}.

\subsection{gem5 Configuration}

\begin{itemize}
    \item \textbf{CPU model:}   TimingSimpleCPU
    \item \textbf{ISA:}         RISC-V RV64GC
    \item \textbf{Caches:}      enabled (default L1 I/D + L2)
    \item \textbf{Memory:}      4\,GB simulated physical RAM
    \item \textbf{Spill mode:}  summary only (\texttt{spill\_verbose=False})
\end{itemize}

The full gem5 invocation was:

\begin{verbatim}
./build/RISCV/gem5.opt
  --outdir=benchmarks/imagick/m5out_test
  configs/deprecated/example/se.py
  --cmd=benchmarks/imagick/imagick_r.riscv
  --options="-limit disk 0
             benchmarks/imagick/test_input.tga
             -shear 25 -resize 640x480
             -negate -alpha Off
             benchmarks/imagick/test_output.tga"
  --cpu-type=TimingSimpleCPU
  --caches
  --mem-size=4GB
\end{verbatim}

% =========================================================
\section{Simulation Results}

The simulation completed successfully in 93 seconds of host wall time
--- the fastest completion among all SPEC CPU2017 benchmarks evaluated
in this study. The output image \texttt{test\_output.tga} was produced
without errors.

\subsection{Pre-ROI Context}

Before the ROI, the binary executed ImageMagick library initialisation:
loading codec registries, configuring memory limits (\texttt{-limit disk 0}),
and parsing command-line options. This pre-ROI phase is extremely small.
Table~\ref{tab:preroi} shows the global counters at ROI entry.

\begin{table}[H]
\centering
\caption{Global Counters at ROI Entry}
\label{tab:preroi}
\begin{tabular}{lr}
\toprule
\textbf{Metric} & \textbf{Value} \\
\midrule
Total instructions (pre-ROI) & 432,667 \\
Total spills detected (pre-ROI) & 14,530 \\
Pre-ROI spill rate & 3.36\% \\
\bottomrule
\end{tabular}
\end{table}

The pre-ROI phase is negligibly small (432K instructions) compared to
all other benchmarks, confirming that virtually all meaningful
computation is inside \texttt{ConvertMain()} and thus inside the ROI.

\subsection{ROI Spill Statistics}

Table~\ref{tab:roi} presents the spill statistics collected within
the ROI.

\begin{table}[H]
\centering
\caption{ROI Spill Statistics --- 538.imagick\_r (test)}
\label{tab:roi}
\begin{tabular}{lr}
\toprule
\textbf{Metric} & \textbf{Value} \\
\midrule
ROI instructions         & 117,706,530 \\
ROI stores               &   8,985,594 \\
ROI loads                &  20,607,956 \\
\midrule
\textbf{ROI spills}      & \textbf{2,424,899} \\
\textbf{Spill rate}      & \textbf{2.0601\%} \\
\bottomrule
\end{tabular}
\end{table}

\subsection{Timing}

\begin{table}[H]
\centering
\caption{Simulation Timing}
\label{tab:timing}
\begin{tabular}{lr}
\toprule
\textbf{Metric} & \textbf{Value} \\
\midrule
Simulated time           & 0.23\,s \\
Host wall time           & 93.10\,s ($\approx$1.6\,min) \\
Total simulated instructions & 118,567,806 \\
\bottomrule
\end{tabular}
\end{table}

% =========================================================
\section{Analysis}

\subsection{Spill Rate Interpretation}

The measured spill rate of \textbf{2.06\%} means that approximately
1 in every 49 instructions inside the ROI is a register spill.
Table~\ref{tab:compare} places imagick\_r in the lower half of the
spill rate distribution across all evaluated benchmarks.

\begin{table}[H]
\centering
\caption{Cross-Benchmark Spill Rate Comparison}
\label{tab:compare}
\begin{tabular}{lrr}
\toprule
\textbf{Benchmark} & \textbf{ROI Spills} & \textbf{Spill Rate} \\
\midrule
507.cactuBSSN\_r & 3,782,099,438 & 7.8756\% \\
500.perlbench\_r &     1,070,055 & 7.2838\% \\
526.blender\_r   &    61,845,869 & 6.2193\% \\
541.leela\_r     & 1,158,236,271 & 4.4438\% \\
519.lbm\_r       &   274,130,440 & 4.2130\% \\
544.nab\_r       &   151,310,616 & 2.2612\% \\
538.imagick\_r   &     2,424,899 & \textbf{2.0601\%} \\
505.mcf\_r       &   446,113,136 & 1.8263\% \\
508.namd\_r      &   327,718,955 & 1.0269\% \\
\bottomrule
\end{tabular}
\end{table}

\subsection{Store vs.\ Load Balance}

The load count (20.6M) is approximately $2.29\times$ higher than the
store count (9.0M). This reflects the read-heavy nature of image
processing pipelines: each pixel is read for computation (potentially
multiple times for filtering kernels), but written back only once.
For the shear and resize operations, source pixels are accessed via
interpolation, reading multiple neighbours per output pixel.

\subsection{Spill Fraction of Loads}

Within the ROI, the ratio of spills to total loads is:
\[
\frac{2{,}424{,}899}{20{,}607{,}956} \approx 11.8\%
\]
Approximately 1 in 8 load instructions is a spill reload. This is
notably lower than the compute-intensive benchmarks (cactuBSSN: 30.3\%,
perlbench: 30.6\%), reflecting ImageMagick's pipeline-oriented
architecture: each transformation stage has a relatively simple
per-pixel computation with a modest number of simultaneously live
variables.

\subsection{ROI Coverage}

The ROI accounts for $117{,}706{,}530$ out of $118{,}567{,}806$
total simulated instructions:
\[
\frac{117{,}706{,}530}{118{,}567{,}806} \approx 99.3\%
\]
This is the highest ROI coverage fraction across all benchmarks
(tied with cactuBSSN's 98.1\%), confirming that the ROI placement
around \texttt{ConvertMain()} captures virtually all benchmark
computation. The 0.7\% outside the ROI (432K instructions) is
exclusively library initialisation.

\subsection{Absolute Instruction Count}

With only 117.7 million ROI instructions, imagick\_r has the smallest
absolute ROI instruction count of any benchmark evaluated. The small
test image and limited number of operations (shear, resize, negate,
alpha removal) result in a workload that is representative in character
but brief in duration. The refrate dataset, with its more expensive
operations (edge detection with radius 41, 181\% resampling, emboss,
mean-shift) on a larger image, would be substantially more intensive.

% =========================================================
\section{Conclusion}

The \texttt{538.imagick\_r} benchmark exhibits a spill rate of 2.06\%
within its image transformation ROI on RISC-V RV64GC. With 2.4 million
spill events across 117 million ROI instructions, imagick\_r has the
smallest absolute spill count of any evaluated benchmark. The
near-complete ROI coverage (99.3\%) confirms the instrumentation
accuracy, while the modest spill rate (1 in 49 instructions) reflects
ImageMagick's pipeline-oriented, per-pixel computation model with
limited simultaneously live variables. At 93 seconds host time,
imagick\_r is by far the fastest benchmark in this study, suggesting
that the test workload may underrepresent the full benchmark character
relative to the refrate dataset.

\end{document}
